\documentclass[11pt,a4paper]{article}

\usepackage[utf8]{inputenc}
\usepackage[T1]{fontenc}
\usepackage[english]{babel}
\usepackage{geometry}
\geometry{margin=2.5cm}
\usepackage{graphicx}
\usepackage{amsmath,amssymb}
\usepackage{booktabs}
\usepackage{caption}
\usepackage{subcaption}
\usepackage{hyperref}
\usepackage{xcolor}
\usepackage{listings}
\usepackage{enumitem}

\hypersetup{
    colorlinks=true,
    linkcolor=blue,
    citecolor=blue,
    urlcolor=blue
}

\lstset{
    basicstyle=\ttfamily\small,
    breaklines=true,
    columns=fullflexible,
    backgroundcolor=\color{gray!10},
    frame=single
}

\newcommand{\todo}[1]{\textcolor{red}{\textbf{[TODO: #1]}}}

\title{Selection Bias and Biomarker Stability in High-Dimensional\\
Leukemia Gene Expression Classification}
\author{\todo{Your Name(s)} \\ \todo{Course / Institution}}
\date{\today}

\begin{document}
\maketitle

\begin{abstract}
We address the problem of classifying leukemia subtypes (Acute LymphoblasticLeukemia, ALL, vs. Acute Myeloblastic Leukemia, AML) from $p=7129$ gene expression probes measured on $n=72$ patients.
Because $p \gg n$, naive model-selection pipelines that perform feature selection \emph{outside} cross-validation produce severely optimistic error estimates, a phenomenon first quantified empirically by Ambroise and McLachlan (2002).
We replicate their experiment by contrasting a \emph{biased} pipeline (feature selection performed once, on the full dataset, before cross-validation) against an \emph{unbiased} pipeline (feature selection re-performed inside every training fold) for two feature-selection methods (SVM-RFE and $L_1$-penalized logistic regression).
We further quantify the \emph{stability} of selected biomarkers via gene inclusion probabilities across folds, contrast them with genes flagged by a univariate FDR-controlled ANOVA test, and report final hold-out performance using only genes that are stable across resampling.
\end{abstract}

\tableofcontents
\newpage

%==================================================================
\section{Introduction}
%==================================================================
High-throughput gene expression assays (microarrays, RNA-seq) allow simultaneous measurement of thousands of transcripts from a single tissue sample.
In precision oncology this is used to distinguish molecularly similar but clinically distinct cancer subtypes -- for example, Acute Lymphoblastic Leukemia (ALL) vs. Acute Myeloblastic Leukemia (AML) -- in order to choose the correct chemotherapy protocol.

Statistically, this is a $p \gg n$ classification problem: thousands of candidate predictors (genes), only a few dozen samples, and strong collinearity between genes that participate in the same biological pathways.
This combination creates two well-known pitfalls that are the core focus of this report:

\begin{enumerate}
    \item \textbf{Selection bias}: if features are selected using the same data later used to estimate the model's generalization error (e.g. selecting genes on the full dataset and \emph{then} cross-validating the classifier), the resulting accuracy/error estimates are systematically optimistic, sometimes dramatically so \cite{ambroise2002selection}.
    \item \textbf{Selection instability}: with collinear, high-dimensional features, the \emph{identity} of the "best" gene subset is highly sensitive to which samples happen to be in the training fold. A model can have good accuracy while selecting a different, largely non-reproducible, set of genes every time it is refit.
\end{enumerate}

Our goals, following the assignment, are therefore not only to produce a classifier with low test error, but to (i) explicitly measure and document the magnitude of the selection bias on this dataset, (ii) report a measure of \emph{biomarker stability} (inclusion probability), and (iii) compare the genes selected by machine-learning wrapper/embedded methods (SVM-RFE, Lasso-logistic-regression) against the genes flagged by a classical univariate test with multiple-testing correction (FDR, Benjamini--Hochberg).

%==================================================================
\section{Data}
%==================================================================
The dataset originates from a clinical leukemia stratification study and contains $n=72$ patients described by:
\begin{itemize}
    \item 5 clinical/demographic factors: \texttt{BM.PB} (bone marrow vs. peripheral blood), \texttt{Gender}, \texttt{Source} (originating hospital), \texttt{tissue.mf} (combination of source and gender), and \texttt{Samples} (identifier);
    \item the target label \texttt{cancer} $\in \{\texttt{allB}, \texttt{allT}, \texttt{aml}\}$;
    \item $7129$ numeric gene-probe expression values.
\end{itemize}

\subsection{Binarization of the target}
The assignment explicitly evaluates models on \emph{binary} classification error, while \texttt{cancer} has three levels.
We therefore collapse the two ALL subtypes into a single class, giving the clinically standard binary target
\[
y \in \{\texttt{ALL}, \texttt{AML}\}, \qquad
\texttt{ALL} = \{\texttt{allB}, \texttt{allT}\}.
\]
This step is performed in \texttt{preprocess.py} \emph{before} the train/test split, and is controlled by a configuration flag (\texttt{binarize\_target}) so that the original 3-class problem remains available for completeness/ablation.

\subsection{Train / test split}
We hold out a dedicated test set (e.g. $80\%/20\%$, stratified on the binarized label) that is \textbf{never touched} during model selection, feature selection, or hyperparameter tuning -- it is only used once, at the very end, by \texttt{evaluate.py}.
Stratified, shuffled splitting is used to avoid degenerate folds, since with $n=72$ samples an unshuffled split could easily produce a test set missing one class entirely.

%==================================================================
\section{Methodology}
%==================================================================

\subsection{Preprocessing}
Clinical/demographic columns and the sample identifier are dropped, leaving only the $7129$ gene expression features as predictors.
Standardization (zero mean, unit variance) is always fit \emph{only} on the current training data (whether that is the full CV set for diagnostic plots, or a single CV fold during model selection) and applied unchanged to the corresponding validation/test partition, to avoid leaking information about the held-out samples into the scaling statistics.

\subsection{Feature selection methods}
Three independent feature-selection criteria are used:

\begin{description}
    \item[Univariate FDR test] A one-way ANOVA $F$-test per gene comparing expression across classes, followed by Benjamini--Hochberg correction for multiple testing at level $\alpha$ (typically $0.05$). This provides a purely statistical baseline that does not use any classifier.
    \item[SVM-RFE] Recursive Feature Elimination with a linear-kernel SVM as the base estimator \cite{guyon2002gene}: at each iteration the $10\%$ least-important features (by SVM weight magnitude) are pruned, until exactly $k$ genes remain.
    \item[Lasso (L1-logistic)] $L_1$-penalized logistic regression (\texttt{solver="saga"}), where a gene is "selected" if its coefficient is non-zero. The inverse regularization strength $C$ controls sparsity directly.
\end{description}

\subsection{Biased vs. unbiased cross-validation}
To replicate the selection-bias experiment of Ambroise and McLachlan \cite{ambroise2002selection}, we compare two pipelines for both SVM-RFE and Lasso:

\begin{itemize}
    \item \textbf{Biased pipeline}: feature selection (SVM-RFE / Lasso) is performed \emph{once}, on the entire CV dataset (i.e. including what will later become the validation folds). A linear SVM is then trained and validated with repeated stratified $k$-fold CV using only the pre-selected genes.
    \item \textbf{Unbiased pipeline}: for \emph{every} training fold, feature selection is re-run from scratch on the training portion only, and the corresponding validation fold is scored using only the genes selected from that fold's training data.
\end{itemize}

Both pipelines are evaluated for several candidate numbers of selected genes $k \in \{10, 20, 30, 50, 100, \dots\}$ (configurable), using $\texttt{n\_splits}\times\texttt{n\_repeats}$ stratified CV.
We report mean training/validation accuracy and log-loss for both pipelines as a function of $k$, which reproduces Figure 2 of \cite{ambroise2002selection}: the biased validation curve is optimistically close to the training curve (small or negative bias-variance gap), while the unbiased validation curve exposes the true, larger generalization gap.

\subsection{Biomarker stability: inclusion probabilities}
Within the unbiased CV loop, every time a gene is selected by SVM-RFE, FDR, or Lasso on a given training fold, a per-gene counter is incremented.
Dividing by the total number of CV iterations ($\texttt{n\_splits}\times\texttt{n\_repeats}$) yields an
\textbf{inclusion probability} per gene, per method.
Genes with high inclusion probability ($>0.5$, say) are considered \emph{stable} biomarkers: their selection does not depend on which 80\% of patients happened to be in the training fold.
This directly targets the assignment's "Variable Selection Stability" requirement and is the natural antidote to the instability induced by collinear gene expression features.

We also report how many genes are jointly stable ($>50\%$ inclusion probability) across \emph{all three} methods (FDR, SVM-RFE, Lasso) -- this addresses the "Control of False Discoveries" requirement, since genes that are both statistically significant \emph{and} repeatedly chosen by two independent ML methods are much less likely to be false discoveries than genes chosen by only one method, by chance, on one particular fold.

\subsection{Final model and held-out evaluation}
The configuration of $k$ (number of genes) that minimizes the unbiased validation log-loss is selected as "best".
Using the inclusion probabilities obtained from the corresponding unbiased CV run, we pick the top-$k$ most stable genes for SVM-RFE and, separately, for Lasso. A linear SVM is then refit on the \emph{entire} CV set restricted to each stable gene subset, and evaluated exactly once on the untouched test set.
This final number is the only one in the report that can be interpreted as an unbiased estimate of real-world generalization performance, since it is the only step where the test set is used.

%==================================================================
\section{Implementation notes and corrections applied}
%==================================================================
During development of the pipeline (\texttt{preprocess.py}, \texttt{feature\_selection.py}, \texttt{cross\_validation.py}, \texttt{bias\_experiment.py}, \texttt{evaluate.py}, \texttt{plotting.py}, \texttt{main.py}), the following issues were identified and fixed to ensure correctness of the final results:

\begin{enumerate}
    \item \textbf{Inconsistent result-table column} used to build the bias-vs-genes plots: the column collecting the number of selected genes must be appended once per configuration (not once per \texttt{(genes, alpha, C)} triplet), and must be named consistently between the experiment runner and the plotting function.
    \item \textbf{Index mismatch} when forwarding RFE/Lasso inclusion probabilities from the unbiased cross-validation routine to the final evaluation step (RFE and FDR probabilities were swapped).
    \item \textbf{Train/test splitting} must use shuffling when stratifying by class, otherwise the split is not well-defined for this small, ordered dataset.
    \item \textbf{Target binarization} (ALL vs AML) is performed explicitly and documented, since the raw label has three levels but the task requires a binary error metric.
    \item \textbf{Multi-class log-loss} calls explicitly pass the fitted classifier's \texttt{classes\_} as labels, to avoid undefined behaviour when a cross-validation fold does not contain every class.
    \item Default hyperparameter grids in the configuration loader are proper lists (e.g. number of genes, FDR $\alpha$, SVM/Lasso $C$), not scalars, since the bias experiment iterates over them.
\end{enumerate}

%==================================================================
\section{Results}
%==================================================================

\subsection{Correlation structure of the raw features}
\todo{Insert \texttt{correlation\_matrix.pdf} (or a representative subset)
here, and comment on the degree of multicollinearity observed among gene
probes.}

\subsection{Selection bias: biased vs. unbiased cross-validation}
\begin{figure}[h!]
    \centering
    \todo{Insert \texttt{bias\_comparison\_accuracy.pdf}}
    \caption{Mean training/validation accuracy for the biased and unbiased
    pipelines, as a function of the number of selected genes, for SVM-RFE
    (left) and Lasso (right).}
    \label{fig:bias-acc}
\end{figure}

\begin{figure}[h!]
    \centering
    \todo{Insert \texttt{bias\_comparison\_loss.pdf}}
    \caption{Mean training/validation log-loss, biased vs. unbiased.}
    \label{fig:bias-loss}
\end{figure}

\todo{Discuss the gap between the biased and unbiased validation curves
in Figures~\ref{fig:bias-acc}--\ref{fig:bias-loss}. State explicitly, with
numbers, how much the biased pipeline overestimates accuracy / underestimates
loss relative to the unbiased pipeline, and at which $k$ the gap is largest.}

\subsection{Biomarker stability}
\begin{figure}[h!]
    \centering
    \todo{Insert \texttt{inclusion\_prob\_histograms.pdf} and
    \texttt{inclusion\_prob\_ranked.pdf}}
    \caption{Distribution and ranked inclusion probabilities for SVM-RFE,
    FDR, and Lasso.}
    \label{fig:stability}
\end{figure}

\begin{table}[h!]
    \centering
    \caption{Top-10 most stable genes under SVM-RFE (highest inclusion
    probability across repeated CV folds).}
    \label{tab:top-genes}
    \begin{tabular}{lc}
        \toprule
        Gene probe & Inclusion probability \\
        \midrule
        \todo{gene 1} & \todo{0.xx} \\
        \todo{gene 2} & \todo{0.xx} \\
        \dots & \dots \\
        \bottomrule
    \end{tabular}
\end{table}

\todo{Report the number of genes jointly stable ($>50\%$) across FDR,
SVM-RFE \emph{and} Lasso, and briefly comment on whether the methods agree
more than chance would predict.}

\subsection{Final held-out performance}
\todo{Report final test-set accuracy (and, if relevant, log-loss / confusion
matrix) for the SVM trained on the top stable RFE genes and on the top
stable Lasso genes, at the best $k$ identified by the unbiased CV log-loss.}

%==================================================================
\section{Discussion}
%==================================================================
The unbiased pipeline is expected to show a markedly larger gap between training and validation performance than the biased pipeline, especially for small $k$ (strong feature selection) where the opportunity to overfit by "shopping" for predictive genes across the whole dataset is greatest.
This is the central, reproducible finding of Ambroise and McLachlan \cite{ambroise2002selection}: with $p \gg n$ and feature selection performed outside the resampling loop, cross-validated error estimates can be arbitrarily optimistic, even on data with no true signal (a sanity check that can be obtained by re-running the same pipeline on permuted labels, where the biased method should still report better-than-chance accuracy due to leakage, while the unbiased method should correctly collapse to the chance level $1/n_{\text{classes}}$).

Stability analysis is expected to show that, while SVM-RFE and Lasso each select on the order of $k$ genes per fold by construction, only a fraction of those have inclusion probability above 50\%, confirming that the \emph{identity} of selected genes is far less reproducible than the resulting predictive accuracy would suggest -- precisely the instability that motivates reporting inclusion probabilities rather than a single "best" gene list.

\todo{Add your own commentary once the experiment has actually been run on
your environment: does the FDR set overlap substantially with the
ML-selected stable genes? Does increasing $k$ improve stability at the cost
of sparsity? Which of SVM-RFE or Lasso produces more stable selections on
this dataset, and why might that be (e.g. Lasso's tendency to arbitrarily
pick one representative from a group of correlated genes vs. RFE's
iterative ranking)?}

%==================================================================
\section{Conclusion}
%==================================================================
We implemented and validated a complete pipeline for high-dimensional gene expression classification that explicitly separates feature selection from performance estimation, quantifies the resulting selection bias by replicating the Ambroise--McLachlan experiment, and reports biomarker stability via per-gene inclusion probabilities under repeated cross-validation.
The methodology -- rather than the raw test accuracy of any single model -- is the primary contribution requested by the assignment: a small, sparse, and cross-validated-stable set of biomarkers is of much greater scientific and clinical value than a marginally higher accuracy obtained through a biased, non-reproducible selection procedure.

%==================================================================
\begin{thebibliography}{9}

\bibitem{ambroise2002selection}
C. Ambroise and G. J. McLachlan,
\emph{Selection bias in gene extraction on the basis of microarray gene-expression data},
Proceedings of the National Academy of Sciences, 99(10), 6562--6566, 2002.

\bibitem{guyon2002gene}
I. Guyon, J. Weston, S. Barnhill, V. Vapnik,
\emph{Gene Selection for Cancer Classification using Support Vector Machines},
Machine Learning, 46, 389--422, 2002.

\end{thebibliography}

\end{document}